%%--------------------------------------------------------------------------
%%  SECTION 3: MATHEMATICAL THEORY OF DEEP NEURAL NETWORK TRAINING
%%--------------------------------------------------------------------------
\section{Mathematical Theory of Deep Neural Network Training}

\noindent\textit{Why it matters.}\enspace
Deploying neural networks in safety-critical scientific computing demands
more than empirical validation: it requires a rigorous mathematical
understanding of training dynamics and generalization.  Without such theory,
one cannot predict when a neural network will fail---an unacceptable
situation in nuclear safety transport, clinical imaging, or autonomous
systems.  Beyond safety, a quantitative theory of learning dynamics is an
essential design tool: one must understand how a network's inductive bias
interacts with the spectral properties of the target PDE solution to
guarantee convergence and inform architecture choices.

\subsection{The Frequency Principle}

A central theme of my theoretical work is the \emph{Frequency Principle}
(F-principle)~\cite{FP1}: the finding, first systematically documented by
our group, that gradient descent on neural network loss surfaces fits
low-frequency components of the target function first and progressively
captures higher frequencies.  This is \emph{a priori} surprising from an
approximation-theoretic standpoint, and its mathematical explanation reveals
a deep connection between the spectral properties of the neural tangent
kernel and the implicit regularization induced by early stopping.

In~\cite{FP1} we provided the first systematic Fourier analysis
establishing F-principle for two-layer networks, with extensive controlled
experiments spanning regression, classification, and PDE residual
minimization.  Subsequent work established:
\begin{itemize}
  \item A \emph{general theory}~\cite{Luo2021} for deep networks of
    arbitrary width and depth, analyzing the spectral dynamics of the
    infinite-width neural tangent kernel and yielding closed-form predictions
    for training convergence rates.

  \item An \emph{exactly solvable linear ODE model}~\cite{Luo2022} that
    reproduces F-principle dynamics for two-layer networks, enabling analytic
    predictions of convergence as a function of architecture and learning rate
    without any approximation.

  \item \emph{Phase diagrams}~\cite{PD21} characterizing the distinct
    generalization phases of two-layer ReLU networks at the infinite-width
    limit, providing a rigorous map from initialization and architecture to
    implicit bias and test error.
\end{itemize}

These theoretical results carry concrete engineering consequences for the
design of scientific machine learning methods.  The F-principle predicts
that na\"{i}ve PINN formulations will fail to resolve high-frequency
components of multiscale PDE solutions in stiff regimes---exactly the failure
mode that motivates the APNN architecture.  Conversely, the exact ODE
model of~\cite{Luo2022} provides quantitative guidance on how network depth
and width should scale with the regularity of the target PDE solution, turning
previously heuristic design choices into mathematically principled decisions.
